\section{All odd rank-two symplectic cores above the small endpoint}
\label{sec:odd-c2-all-fields}

\paragraph{Status and exact scope.}
This complete derivation and its cited input scopes have passed whole
author, root and independent peer readings. The corrected unipotent
input and finite spectra have separate exact certificates.
The core factors in the final theorem are permuted transitively;
this is the normal-core configuration furnished by D-BI-19 for the
primitive, factorwise tensor-indecomposable source gate. No assertion
about an arbitrary intransitive mixture is substituted for that
hypothesis. The original inducing problem has a nonselfdual complex
irreducible representation, in particular one with exact imaginary
quadratic character field. Its inducing field may be larger.

Write $\nu(T)=a-\max_\lambda\dim\ker(T-\lambda I)$ for projective
support on an $a$-dimensional space. Scalar multiplication preserves it,
ordinary moving rank is at least it, and $\nu(T)\ge\nu(T^j)$.
Every nontrivial irreducible representation of $\operatorname{Sp}_4(q)$,
$q\ge5$, has projective kernel exactly its center: that kernel is
normal, its image in the simple central quotient is either trivial
or the whole quotient, and the latter would make a perfect group
act by scalars and hence trivially.

\begin{lemma}[D-OC2-01: the corrected unipotent input]
\label{D-OC2-01}
For odd $q\ge7$, every nontrivial ordinary irreducible character $\chi$
of $S=\operatorname{Sp}_4(q)$, of degree $a$, satisfies
\[
 a\ge(q^2-1)/2,\qquad |\chi(u)|\le a/2
       \quad(1\ne u\in S\text{ unipotent}).                 \tag{OC2.1}
\]
Apart from the four characters of degrees $(q^2+1)/2$ and
$(q^2-1)/2$, every such degree is at least $q(q-1)^2/2$.
For $q\equiv3\pmod4$ those four characters are nonselfdual.
\end{lemma}

\begin{proof}
We use the published corrected table of values of \emph{all ordinary
characters on unipotent elements} in
\cite[pp.~620--624, Table~2 and the restriction formulas]{OC2:Dabbaghian2006}.
This is a different input from an incomplete generic full table or a
table of only unipotent characters. There are six nonidentity
unipotent classes, denoted $A_{21},A_{22},A_{31},A_{32},A_{41},A_{42}$,
and $43$ nontrivial character types. Some parameterized types can be
empty at particular small $q$, which does not impair an upper bound
checked on every type. The degree column and these six values are
retained in \texttt{unipotent\_data.json}.

Here are the mathematical inputs and the exact coefficient procedure.
Put $t^2=(-1)^{(q-1)/2}q$ and $\epsilon=(-1+t)/2$,
$\epsilon'=(-1-t)/2$. Swapping the sign of $t$ exchanges paired
character labels. The six basis characters in the published
restriction formulas are $a_1,a_2,b,c,d_1,d_2$, with $a_2,d_2$
obtained from $a_1,d_1$ by that exchange. In the column order
$(1,A_{21},A_{22},A_{31},A_{32},A_{41},A_{42})$, the other four rows are
\begin{align*}
a_1={}&\bigl((q^2+1)/2,(q+1)/2+q\epsilon,
 (q+1)/2+q\epsilon',(q+1)/2,(1-q)/2,-\epsilon',-\epsilon\bigr),\\
b={}&\bigl(q(q-1)^2/2,q(1-q)/2,q(1-q)/2,0,q,0,0\bigr),\\
c={}&\bigl(q(q^2+1)/2,q(q+1)/2,q(q+1)/2,0,q,0,0\bigr),\\
d_1={}&\bigl((q-1)^2(q^2+1)/2,
 (1-q)/2-q(q-1)\epsilon',
 (1-q)/2-q(q-1)\epsilon,
 (1-q)/2,(1-3q)/2,-\epsilon,-\epsilon'\bigr).
                                                               \tag{OC2.2}
\end{align*}
The complete integer coefficient vectors in this six-character basis
and the trivial character are retained. Four rows are instead read
\emph{directly from Table~2}, as follows, with each conjugate obtained
by exchanging $t$ and $-t$:
\begin{align*}
\xi_{41}={}&\bigl((q+1)^2(q^2+1)/2,
 (q+1)/2-q(q+1)\epsilon,(q+1)/2-q(q+1)\epsilon',
 (1+3q)/2,(q+1)/2,-\epsilon,-\epsilon'\bigr),\\
-\xi'_{41}={}&\bigl((q^4-1)/2,
 -(1-q)/2+q(q+1)\epsilon,-(1-q)/2+q(q+1)\epsilon',
 -(1-q)/2,-(1+q)/2,\epsilon,\epsilon'\bigr).
                                                               \tag{OC2.3}
\end{align*}
The indexed text of the four corresponding restriction identities
interchanges $d_1,d_2$ relative to Table~2. We do not use those
inconsistent identities. The initial failed comparison and the direct
Table~2 overrides are retained explicitly; no correction is inferred
merely by fitting finitely many values of $q$.

Every resulting entry has the form $A(q)+B(q)t$ with rational
polynomials $A,B$. For $q\ge7$,
\[
        |t|=\sqrt q\le(q-1)/2,                                \tag{OC2.4}
\]
because $(q-1)^2-4q$ has positive coefficients after $q=7+x$.
Expand $A(7+x),B(7+x)$ and bound each by the polynomial obtained by
taking the absolute values of its coefficients. Multiply the second
bound by $3+x/2$. For all $43\cdot6=258$ entries, subtracting this
upper bound from $a(7+x)/2$ gives a polynomial with nonnegative
rational coefficients. The same exact calculation checks the degree
floor $(q^2-1)/2$, and the stronger floor $q(q-1)^2/2$ away from the
four specified types. The complete coefficient lists and the safe
arithmetic program are retained in \texttt{unipotent\_checked.json}
and \texttt{check\_unipotent.py}. This proves the inequalities for
all $q\ge7$, not just for sampled fields.

As a separate transcription check, the published formulas agree on
all six unipotent classes with all $60$ rows of the complete
$\operatorname{Sp}_4(5)$ table, all $52$ rows of the complete
$\operatorname{PSp}_4(7)$ table, and all $74$ rows of the complete
$\operatorname{PSp}_4(9)$ table, using a single column permutation
for each table. This is an executed exact cyclotomic comparison;
the two latter comparisons do not purport to supply missing
center-negative tables. The source's averaging identities for all six
basis characters provide another exact check on the positions of the
zero entries in (OC2.2).

Finally the four small characters are $a_1,a_2,a_1-1,a_2-1$ on the
unipotent subgroup. When $q\equiv3\pmod4$, their $A_{21}$ values
have a nonzero imaginary coefficient of $\sqrt{-q}$, so none is
selfdual. The four small characters are the four usual symplectic
Weil characters: the two nontrivial additive-character square classes
each give the two parity constituents of degrees $(q^2\pm1)/2$,
as in the already checked D-SPW-01, and this degree list contains
exactly those four characters.
\end{proof}

\begin{lemma}[D-OC2-02: every noncentral inner operator]
\label{D-OC2-02}
For $q\ge7$ odd and $\rho$ as above, every noncentral element
$x\in\operatorname{Sp}_4(q)$ satisfies
\[
                   \nu(\rho(x))\ge \frac{q-1}{4q}\,a.
                                                               \tag{OC2.5}
\]
\end{lemma}

\begin{proof}
Take a power $h$ whose image modulo $\{\pm I\}$ has prime order
$\ell$. Its projective support is no greater than that of $x$.
If $\ell=p$, multiplication by a central sign makes $h$ a nonidentity
unipotent, and (OC2.1) gives support at least $a/4$. Suppose
$\ell\ne p$. Then $h$ is semisimple. If $\ell$ is odd, multiply
by a central sign so that $h$ has actual order $\ell$. If $\ell=2$,
then $h^2=I$ or $h^2=-I$.

We first describe a useful invariant isotropic-subspace case. Let
$R$ be a nonzero totally isotropic $h$-stable natural subspace and
suppose conjugation by $h$ acts nontrivially on the symmetric shear
group attached to $R$. For a basis $r_1,\ldots,r_d$ of $R$ and a
symmetric $d\times d$ matrix $B$, put
\[
 N_B(v)=\sum_{i,j}b_{ij}\langle r_j,v\rangle r_i,
                    u_B=I+N_B.                                \tag{OC2.6}
\]
Symmetry and isotropy give $N_BN_C=0$ and
$\langle N_Bv,w\rangle+\langle v,N_Bw\rangle=0$.
Consequently $u_B$ is symplectic and $u_Bu_C=u_{B+C}$.
If $A=h|_R$ in this basis, conjugation acts by $B\mapsto ABA^t$.
The map $B\mapsto N_B$ is injective, because the pairings with
the independent $r_j$ are independent linear forms.

This conjugation action has order $\ell$: the central square of
$h$, if any, acts trivially on symmetric matrices. Since
$\ell\ne p$, its nonzero image under $h-1$ is an $\mathbf F_q$
subspace on which $h$ has no fixed vectors. It supplies an abelian
unipotent subgroup $U$ of order $Q\ge q$. Each nontrivial dual
weight has an orbit of length $\ell$. The invariant weight has
dimension at most $a/2+a/(2Q)$ by averaging (OC2.1), and every
other orbit contributes at most $1/\ell$ of its dimension to any
chosen eigenspace of $\rho(h)$. This is also valid for a cyclic
block operator with nontrivial block product. Hence
\[
 \nu(\rho(h))\ge(1-1/\ell)(1-1/Q)a/2
                    \ge(q-1)a/(4q).                          \tag{OC2.7}
\]

Both projective involution cases admit such an $R$. If $h^2=I$
and $h\ne\pm I$, its plus and minus eigenspaces are nonzero
orthogonal symplectic planes. Choose a line in each; their sum is
an isotropic plane, and its mixed symmetric shear has eigenvalue
$-1$. If $h^2=-I$ and $-1$ is a square, either eigenplane is
isotropic and its symmetric shears have eigenvalue $-1$. If
$-1$ is not a square, regard the natural space as dimension two
over $E=\mathbf F_{q^2}$, with $h$ multiplication by a square root
of $-1$. The alternating form induces, through the nondegenerate
trace pairing, a nondegenerate hermitian or skew-hermitian form
over $E/\mathbf F_q$. Scaling makes it hermitian. Every such
two-dimensional form over a finite field is isotropic: after
diagonalization this follows by solving one scalar norm equation,
using surjectivity of $E^*\to\mathbf F_q^*$. Its isotropic $E$-line
is an isotropic $\mathbf F_q$-plane $R$. On its symmetric square,
$h$ has a nontrivial minus-one eigenspace, as is seen over $E$
from eigenvalues $i^2=-1$ and $(-i)^2=-1$.

Now let $\ell$ be odd. On a one- or two-dimensional space, a
semisimple $A$ acts trivially on the symmetric square only if
$A=\pm I$: over an algebraic closure its squared eigenvalues
and pairwise products must all equal one. Thus any invariant
isotropic $R$ on which $h$ acts nontrivially is handled by
(OC2.6)--(OC2.7), since $-I$ cannot have odd order.

If an invariant isotropic subspace exists but $h$ fixes it
pointwise, choose a fixed vector $e$ and a fixed paired vector
$f$ with $\langle e,f\rangle=1$; average a paired vector over
the order-$\ell$ group to obtain $f$. The orthogonal complement
$W=\langle e,f\rangle^\perp$ is a symplectic plane. The action
of $h$ on it is nontrivial, and has no fixed vectors; otherwise
its semisimple determinant-one action would be the identity.
The line-parabolic radical consists of the maps
\[
 e\mapsto e,\qquad
 w\mapsto w+\langle v,w\rangle e,\qquad
 f\mapsto f+v+2ze,
             \quad v\in W,\ z\in\mathbf F_q.                  \tag{OC2.8}
\]
Its multiplication is $(v,z)(v',z')=(v+v',z+z'+
\langle v,v'\rangle/2)$, and $h$ acts by $(v,z)\mapsto(hv,z)$.
Every nonidentity element is unipotent. Applying the fully proved
Heisenberg carrier lemma D-HOS-01, with $m=1$, gives
\[
 \nu(\rho(h))\ge\frac23(1-q^{-1})(1-q^{-3})a/2
                       \ge(q-1)a/(4q),                        \tag{OC2.9}
\]
since $1-q^{-3}\ge3/4$.

It remains that $h$ has no invariant isotropic subspace. Its
semisimple primary decomposition has neither a plus/minus-one
summand nor a pair of distinct reciprocal factors, since these
provide an invariant isotropic subspace. A self-reciprocal
irreducible factor of multiplicity at least two also provides
one: its associated hermitian space over the finite extension
has an isotropic line, by the norm argument just given.
Consequently the natural space consists of one self-reciprocal
factor of degree four, or two distinct self-reciprocal factors
of degree two, each with multiplicity one. Their symplectic
centralizers have respective orders $q^2+1$ and $(q+1)^2$.
Indeed their commuting natural operators are field scalars,
and preserving the form imposes the norm-one condition on
each field. Column orthogonality therefore gives
\[
 |\chi(h)|\le\sqrt{|C_S(h)|}\le q+1,
 \qquad \frac{|\chi(h)|}{a}\le\frac2{q-1}.
                                                               \tag{OC2.10}
\]
Thus $\nu\ge a(q-3)/(2(q-1))\ge a(q-1)/(4q)$ for $q\ge7$;
the latter inequality is $q^2-4q-1\ge0$. All cases are exhausted.
\end{proof}

\begin{lemma}[D-OC2-03: compatible outer operators]
\label{D-OC2-03}
For $q\ge7$ odd, every finite-order nonscalar operator normalizing
the represented image of $\rho$ satisfies
\[
 \nu(T)\ge\frac{q-1}{8q}a.                                   \tag{OC2.11}
\]
If $\rho$ is nonselfdual, the stronger bound
\[
 \nu(T)\ge\frac{q-1}{4q}a                                   \tag{OC2.12}
\]
holds, including every compatible field automorphism.
\end{lemma}

\begin{proof}
For (OC2.11), choose a core element with nonidentity commutator
with $T$, using irreducibility. A noncentral commutator has
support at least $(q-1)a/(4q)$, while a nonidentity scalar
commutator has actual moving rank $a$. The elementary rank
inequality $\operatorname{rank}([T,h]-I)\le2\nu(T)$ proves
the claim. It is valid after scalar-normalizing $T$ so that
its largest eigenspace has eigenvalue one.

For the stronger statement we use the exact ordinary automorphism
description $\operatorname{Out}(\operatorname{PSp}_4(q))
=C_2\times C_f$, $q=p^f$, from
\cite[\S2.2.2, Table~2.1, author p.~14]{HES:LiXia2022}.
For $q\ge5$ the universal perfect cover is $\operatorname{Sp}_4(q)$,
with center of order two and no exceptional multiplier; the precise
exception list is \cite[Chapter~5, Table~1, p.~35]{BI:HHM}.
Automorphisms of a represented central quotient lift uniquely to
this cover, and an automorphism trivial on its simple quotient is
trivial: its difference from the identity is a homomorphism from
a perfect group to its center.

Conjugation by a fixed multiplier-minus-one matrix sends every
symplectic conjugacy class to its inverse class. To see the exact
published input, \cite[Theorem~1, author p.~1]{OC2:Vinroot2004}
states that every symplectic element is a product of two
skew-symplectic involutions. If $g=ab$, then $aga=g^{-1}$.
All multiplier-minus-one conjugations differ by an inner
symplectic conjugation, proving the assertion. If $-1$ is a
square, scalar normalization makes this conjugation inner, so
every irreducible character is selfdual. Hence nonselfduality
implies $q\equiv3\pmod4$, $p\equiv3\pmod4$ and $f$ odd.
The nontrivial diagonal outer automorphism $\delta$ then sends
$\rho$ to $\rho^\vee$ and is incompatible with $\rho$.

In fact every compatible outer class has trivial diagonal
component. If $\delta\sigma^j$ stabilized $\rho$, raising it to
the odd order of $\sigma^j$ in $C_f$ would make $\delta$
stabilize $\rho$, a contradiction. A compatible operator is
therefore scalar times an inner operator, or has a nontrivial
field component of odd order. In the latter case take a power
whose field component has prime order $\ell\ge3$; support
cannot increase on taking this power. Write the resulting
automorphism as $\alpha$, with field part $\sigma$ of order
$\ell$. If $\alpha^\ell$ is nontrivial inner, the corresponding
power of the intertwiner is scalar times a noncentral inner
operator, and (OC2.5) proves the desired bound.

Suppose instead $\alpha^\ell=1$. Because its diagonal class is
trivial, it has a natural lift $s=h\sigma$ with
$h\in\operatorname{Sp}_4(q)$. Its $\ell$th power is a central
sign $cI$. Replacing $s$ by $cs$, which implements the same
automorphism, makes $s^\ell=1$, since $\ell$ is odd. Let
$k=\mathbf F_{q_0}$ be the fixed field of $\sigma$, so
$q=q_0^\ell$ and $q_0\ge3$.

We give descent without assuming that $\ell$ is invertible in
the field. For $z\in\mathbf F_q$ and a natural vector $v$, put
\[
             P_z(v)=\sum_{j=0}^{\ell-1}\sigma^j(z)s^j(v).
                                                               \tag{OC2.13}
\]
It is fixed by $s$. Choose a $k$-basis $z_1,\ldots,z_\ell$
of $\mathbf F_q$. The matrix $(\sigma^j(z_i))_{i,j}$ is
invertible. Indeed, a dependence of its columns would give
a nonzero $q_0$-linearized polynomial of degree at most
$q_0^{\ell-1}$ vanishing on a $k$-basis and hence on all
$q_0^\ell$ elements, an impossibility. Reordering the columns
handles any generator of the cyclic Galois group. Consequently
every $v$ is in the $\mathbf F_q$-span of the fixed vectors
$P_{z_i}(v)$. The fixed space has dimension four over $k$:
if $w_i$ are $k$-independent fixed vectors and
$\sum a_iw_i=0$, application of $P_z$ gives
$\sum\operatorname{Tr}_{q/q_0}(za_i)w_i=0$ for every $z$.
Nondegeneracy of the finite-field trace pairing makes every
$a_i=0$. This trace pairing is nondegenerate because its
linearized trace polynomial has degree less than $q$ and
is nonzero, and multiplication by any nonzero scalar is
bijective. No averaging by $1/\ell$ has been used.

The natural alternating form takes values in $k$ on this
fixed space and is nondegenerate there, because the fixed
space spans. A symplectic $k$-basis gives
$P\in\operatorname{Sp}_4(q)$ with $P^{-1}sP=\sigma$.
Relabeling the representation by this inner change, the
intertwiner implements pure Frobenius. Its $\ell$th power is
scalar, so normalize that power to one.

Inside the standard $\operatorname{GL}_2(q)$ Levi take its
one-row additive subgroup $U\simeq(\mathbf F_q,+)$.
The diagonal torus is transitive on its nonzero additive
characters. If each nonzero weight has multiplicity $m$,
averaging (OC2.1) yields $m\ge a/(2q)$. Frobenius fixes
exactly $q_0$ of the $q$ dual weights and puts every other
weight in an orbit of length $\ell$. Each such orbit
contributes at least a fraction $1-1/\ell$ of its dimension
to projective support. Therefore
\[
 \nu(T)\ge\frac12(1-1/\ell)(1-q_0/q)a
       \ge\frac8{27}a>\frac14a
       >\frac{q-1}{4q}a.                                    \tag{OC2.14}
\]
Here $\ell\ge3$ and $q_0/q=q_0^{1-\ell}\le1/9$.
This completes the all-field normalizer proof.
\end{proof}

\begin{theorem}[D-OC2-04: the all-odd-$C_2$ source exclusion]
\label{D-OC2-04}
Let $J\subset\operatorname{GL}(V)$ be a finite group acting irreducibly
on a complex space $V$, and suppose an
irreducibly acting normal central product of quasisimple groups
with simple quotient $\operatorname{PSp}_4(q)$, $q\ge5$ odd,
has its factors permuted transitively by $J$. Suppose $V$ is an
inducing block of a nonselfdual original irreducible representation
in the sixfold signature-$(3,3)$ curve-cover problem. Then the
actual period image has dimension at most eight. All compatible
scalars, outer automorphisms and factor permutations are allowed.
The assertion imposes no polarization determinant condition and
does not equate the inducing and original character fields.
\end{theorem}

\begin{proof}
For $q=5$, D-OC2-F1--F2 already proves the stronger quarter-bound
exclusion. Assume $q\ge7$. Write $V=\bigotimes_{i=1}^k V_i$,
where the transitive core action makes the factor degrees all
equal to $a\ge(q^2-1)/2$. Their selfduality status is also the
same, since automorphisms and isomorphisms preserve it.
Write $e=a^k$. Factor-preserving operators have convolved
normalized eigenspace distributions. A nontrivial factor cycle
has support at least $e/4$, by the weighted basis-word orbit
argument: at most $a$ words are constant, and every other orbit
has length at least two. These statements permit arbitrary
outer intertwiners on the factors and arbitrary multiplicity
of eigenvalues. Nonidentity scalar inertias move all $e$
dimensions.

The actual intermediate-cover factorization gives
$6=2e(h-1)+D$, where $D$ is the sum of genuine inertia moving
ranks. Since $e\ge24$, genus at least two is impossible; genus
one has at most six genuine branch points and actual period
rank at most six. Rank nine would therefore require
\[
                    h=0,\quad R\ge12,\quad D=2e+6.
                                                               \tag{OC2.15}
\]
If the factors are nonselfdual, (OC2.12) and the tensor argument
give moving rank at least $(q-1)e/(4q)$ at every genuine branch.
Consequently
\[
 2e+6\ge12\frac{q-1}{4q}e,
       \qquad e\le\frac{6q}{q-3}\le\frac{21}2,
                                                               \tag{OC2.16}
\]
contradicting $e\ge24$.

Suppose the factors are selfdual. Their tensor product has a
unique invariant nondegenerate bilinear form up to scalar.
Every element of $J$ preserves its line and hence has a
similitude multiplier $\lambda:J\to\mathbf C^*$.
This multiplier is nontrivial. Otherwise $V$ would be selfdual
as a $J$-representation, and its induction to the original
group would also be selfdual, contrary to the hypothesis.
The product-one generating inertia tuple on the genus-zero
intermediate base contains at least two entries with
$\lambda\ne1$. For each, its fixed subspace is totally
isotropic for the bilinear form and has dimension at most
$e/2$. The two entries cost at least $e$ moving dimensions;
the other ten entries cost at least $(q-1)e/(8q)$ each by
(OC2.11) and the tensor argument. Thus
\[
 2e+6\ge e+10\frac{q-1}{8q}e,
                    \qquad e\le\frac{24q}{q-5}.              \tag{OC2.17}
\]
For $q\ge11$, the degree floor $(q^2-1)/2$ is already
strictly greater than this bound: the difference after
multiplication by $2(q-5)$ is
$q^3-5q^2-49q+5$, whose coefficients at $q=11+x$ are
positive. For $q=7,9$, two factors give respectively
$e\ge24^2,40^2$, whereas (OC2.17) gives $e\le84,54$.
There is therefore only one factor.

At $q=7$, selfduality excludes the four small characters by
D-OC2-01, leaving $a\ge126>84$. At $q=9$, the stronger
degree floor is $288>54$, except for the four Weil modules
of degrees $40,41$. Their full-normalizer source exclusion
is precisely the already reviewed D-SPW-02--03; it includes
arbitrary scalar and field operators and requires no
polarization determinant assumption. These are all the
possibilities, completing the contradiction to rank nine.
\end{proof}

\paragraph{What is not inferred.}
The result excludes the stated transitive ordinary symplectic
core configuration. It does not construct a source, does not
infer dominance from a parameter count, and does not extend
the complete finite-table coverage beyond the groups explicitly
listed in D-OC2-F. The $q=3$ endpoint retains its previously
proved separate reduction to the known degree-four and
degree-five possibilities.

\paragraph{Executed new certificates.}
The separate program \texttt{check\_all\_fields.py} verifies the
universal polynomial inequalities used in (OC2.9), (OC2.10),
(OC2.16), and (OC2.17), together with the symbolic symplectic
shear and Heisenberg multiplication identities. Its finite
$\mathbf F_{27}$ descent witness has Galois order equal to the
characteristic: it checks all $27^4=531441$ natural vectors,
finds exactly $3^4=81$ fixed vectors for an explicit symplectic
semilinear cocycle, checks the Moore matrix and trace pairing,
and obtains three fixed dual weights and eight length-three
orbits. This supports the implementation and the delicate
characteristic-three signs; the preceding universal proof
establishes descent in every field. This finite witness is a
separate check. The retained output is
\texttt{all\_fields\_checked.json}.
